\subsection{Initial Data Curation Approach}
\label{sec:data-curation-initial}

In the initial phase of our project, our benchmark dataset curation process also followed 2 stages similar to those described in Section \ref{sec:data-preparation}. However, instead of using MinerU for data extraction to Markdown format and LangGraph for building an agentic loop to generate the QA sets, we followed a more manual and more inefficient approach. 
For data extraction, we manually converted the PDF case reports to markdown using PyMuPDF. However, this approach often led to formatting issues as well as tables and images being misplaced or lost, which required additional manual intervention to correct. 
This limitation motivated us to research other approaches for handling this task much more effectively, leading us to MinerU. This tool significantly improved the quality of our extracted data and reduced the time spent on manual corrections. Regarding the generation of the QA sets, we initially employed a more straightforward prompting method using \textit{gemini-2.5-flash} without an agentic loop. This approach, however, often resulted in incomplete or less coherent question and reasoning pairs, necessitating multiple rounds of manual review and editing. 
After exploring alternative methods, we discovered LangGraph, the framework that allowed us to generate the QA sets with structured output. This method not only improved the coherence and completeness of the generated content but also streamlined the overall process, reducing the need for extensive manual intervention.
